% ═══════════════════════════════════════════════════════════════════════════
% LEHRERHANDREICHUNG DS-04: Letzte Bucharbeit vor Assessment
% Kompakter 1-2-Seiter mit allem Wichtigen
% ═══════════════════════════════════════════════════════════════════════════

\documentclass[11pt, a4paper]{article}

\usepackage[utf8]{inputenc}
\usepackage[T1]{fontenc}
\usepackage[ngerman]{babel}
\usepackage{helvet}
\renewcommand{\familydefault}{\sfdefault}

\usepackage[margin=1.2cm, top=1cm, bottom=1cm]{geometry}
\usepackage[table]{xcolor}
\usepackage{array}
\usepackage{tabularx}
\usepackage{booktabs}
\usepackage{tikz}
\usepackage{tcolorbox}
\usepackage{enumitem}
\usepackage{amssymb}
\tcbuselibrary{skins}

% Farben
\definecolor{spanisch}{HTML}{c41e3a}
\definecolor{gruppeA}{HTML}{2a7a4b}
\definecolor{gruppeB}{HTML}{1565c0}
\definecolor{hellgrau}{HTML}{f0f0f0}
\definecolor{phase}{HTML}{666666}
\definecolor{wichtig}{HTML}{b35c00}

\pagestyle{empty}
\setlist{nosep, leftmargin=1.5em}

\begin{document}

% ═══════════════════════════════════════════════════════════════════════════
% HEADER
% ═══════════════════════════════════════════════════════════════════════════

\begin{tikzpicture}[remember picture, overlay]
\fill[spanisch] ([yshift=-1cm]current page.north west) rectangle ([yshift=-2.2cm]current page.north east);
\node[white, font=\Large\bfseries, anchor=west] at ([yshift=-1.6cm, xshift=1.2cm]current page.north west) {DS 04 -- Letzte Bucharbeit + Kriterien};
\node[white, font=\normalsize, anchor=east] at ([yshift=-1.6cm, xshift=-1.2cm]current page.north east) {90 Min | 13.01.2026};
\end{tikzpicture}

\vspace{1.2cm}

% ═══════════════════════════════════════════════════════════════════════════
% ÜBERSICHT
% ═══════════════════════════════════════════════════════════════════════════

\begin{minipage}[t]{0.48\textwidth}
\textbf{\textcolor{gruppeA}{GRUPPE A (7 SuS) -- A2}}\\[0.2cm]
\begin{tabular}{@{}ll@{}}
\textbf{Thema:} & Punto final + Evaluación \\
\textbf{Buch:} & S. 33--34 \\
\textbf{Aufgaben:} & Punto final, Evaluación 1 \\
\end{tabular}

\vspace{0.2cm}

\textbf{Lernziele:}
\begin{itemize}
\item E-Mail an Austauschschüler schreiben
\item Evaluación 1 zur Selbstkontrolle
\item Bewertungskriterien verstehen
\item Partner für Assessment bilden
\end{itemize}
\end{minipage}
\hfill
\begin{minipage}[t]{0.48\textwidth}
\textbf{\textcolor{gruppeB}{GRUPPE B (3 SuS) -- B1}}\\[0.2cm]
\begin{tabular}{@{}ll@{}}
\textbf{Thema:} & Vida laboral + Vortrag \\
\textbf{Buch:} & S. 106--107 \\
\textbf{Aufgaben:} & 3b, 4, 5 \\
\end{tabular}

\vspace{0.2cm}

\textbf{Lernziele:}
\begin{itemize}
\item Arbeitsleben DE/ES vergleichen
\item Wortschatzarbeit (Synonyme)
\item Vortrags-Gliederung erstellen
\item Bewertungskriterien verstehen
\end{itemize}
\end{minipage}

\vspace{0.2cm}

\begin{tcolorbox}[
    colback=wichtig!15,
    colframe=wichtig,
    boxrule=1pt,
    top=2pt, bottom=2pt, left=4pt, right=4pt
]
\footnotesize
\textbf{ASSESSMENT am Mo 19.01.2026 (DS 5)!} Heute: Letzte Übung + Kriterien besprechen + Partnerbildung (A) / Themenwahl (B).
\end{tcolorbox}

\vspace{0.2cm}
\hrule
\vspace{0.2cm}

% ═══════════════════════════════════════════════════════════════════════════
% VERLAUFSPLAN
% ═══════════════════════════════════════════════════════════════════════════

\textbf{\large Verlaufsplan}

\vspace{0.2cm}

\footnotesize
\begin{tabularx}{\textwidth}{|c|c|X|X|c|}
\hline
\rowcolor{spanisch!20}
\textbf{Zeit} & \textbf{Phase} & \textbf{Gruppe A} & \textbf{Gruppe B} & \textbf{L bei} \\
\hline
\multicolumn{5}{|c|}{\cellcolor{hellgrau}\textbf{1. STUNDE}} \\
\hline
0--5 & Einstieg & \multicolumn{2}{c|}{``Was steht nächsten Montag an?'' Assessment wiederholen} & alle \\
\hline
5--10 & Überleitung & Punto final erklären (S. 33) & Aufg. 3b--4 erklären (S. 106) & alle \\
\hline
10--40 & Arbeitsphase & \textbf{Punto final:} E-Mail schreiben (EA/PA) & \textbf{Aufg. 3b--4:} Textarbeit (EA) & wechselnd \\
\hline
40--45 & Sicherung & \multicolumn{2}{c|}{1--2 Ergebnisse kurz vorstellen} & alle \\
\hline
\multicolumn{5}{|c|}{\cellcolor{hellgrau}\textbf{2. STUNDE}} \\
\hline
0--20 & Arbeitsphase & \textbf{Evaluación 1} (EA, Selbsttest) & \textbf{Aufg. 5 + Vortrags-Gliederung} & wechselnd \\
\hline
\rowcolor{wichtig!20}
20--35 & \textbf{Kriterien} & \multicolumn{2}{c|}{\textbf{Bewertungskriterien besprechen (getrennt, s.u.)}} & \textbf{beide} \\
\hline
35--42 & Vorbereitung & \textbf{Partnerbildung} für Assessment & \textbf{Thema wählen} für Vortrag & beide \\
\hline
42--45 & Abschluss & \multicolumn{2}{c|}{HA aufgeben, Mut zusprechen!} & alle \\
\hline
\end{tabularx}

\normalsize
\vspace{0.2cm}

% ═══════════════════════════════════════════════════════════════════════════
% MATERIAL & SCAFFOLDS
% ═══════════════════════════════════════════════════════════════════════════

\begin{minipage}[t]{0.48\textwidth}
\begin{tcolorbox}[
    colback=hellgrau,
    colframe=phase,
    title={\footnotesize\bfseries Material-Checkliste},
    fonttitle=\bfseries,
    boxrule=0.5pt,
    top=2pt, bottom=2pt, left=4pt, right=4pt
]
\footnotesize
\begin{itemize}
\item[$\square$] Lehrwerk A\_tope.com
\item[$\square$] \textbf{MATERIAL-05-Bewertungsbögen} (Wiederholung)
\item[$\square$] \textbf{LOESUNG-04-evaluacion1} (Selbstkontrolle A)
\item[$\square$] \textbf{SCAFFOLD-04-beispiel-email} (für A)
\item[$\square$] \textbf{SCAFFOLD-04-satzanfaenge} (für H.P., A.A., A.S.)
\end{itemize}
\end{tcolorbox}
\end{minipage}
\hfill
\begin{minipage}[t]{0.48\textwidth}
\begin{tcolorbox}[
    colback=spanisch!10,
    colframe=spanisch,
    title={\footnotesize\bfseries Hausaufgabe},
    fonttitle=\bfseries,
    boxrule=0.5pt,
    top=2pt, bottom=2pt, left=4pt, right=4pt
]
\footnotesize
\textbf{Gruppe A:}
\begin{itemize}
\item E-Mail fertigstellen (falls nicht fertig)
\item \textbf{Assessment üben:} Mit Partner Fragen üben!
\end{itemize}
\textbf{Gruppe B:}
\begin{itemize}
\item Vortrags-Gliederung fertigstellen
\item \textbf{Assessment üben:} Vortrag laut üben!
\end{itemize}
\end{tcolorbox}
\end{minipage}

\vspace{0.2cm}

% ═══════════════════════════════════════════════════════════════════════════
% ARBEITSBLÄTTER
% ═══════════════════════════════════════════════════════════════════════════

\begin{tcolorbox}[
    colback=hellgrau,
    colframe=phase,
    title={\footnotesize\bfseries Arbeitsblätter \& Scaffolds},
    fonttitle=\bfseries,
    boxrule=0.5pt,
    top=2pt, bottom=2pt, left=4pt, right=4pt
]
\footnotesize
\begin{tabularx}{\textwidth}{@{}lp{3cm}lX@{}}
\textbf{AB-04-A} & Punto final (A) & \textit{optional} & Strukturierte E-Mail-Vorlage. \textit{Einsatz: Alternative zum Buch.} \\[0.2em]
\textbf{AB-04-B} & Vida laboral (B) & \textit{optional} & Erweiterte Textarbeit. \textit{Einsatz: Vertiefung Aufg. 3b--4.} \\[0.2em]
\hline
\textbf{SCAFFOLD} & Beispiel-E-Mail & \textbf{empfohlen} & Muster-E-Mail für Punto final. \textit{Zeigen bei Unsicherheit!} \\[0.2em]
\textbf{SCAFFOLD} & Satzanfänge & \textbf{für H.P., A.A., A.S.} & Formulierungshilfen hay/está/es. \\[0.2em]
\textbf{LÖSUNG} & Evaluación 1 & \textbf{für alle A} & Selbstkontrolle mit Selbsteinschätzung. \\
\end{tabularx}
\end{tcolorbox}

\vspace{0.2cm}

% ═══════════════════════════════════════════════════════════════════════════
% KRITERIEN-BESPRECHUNG
% ═══════════════════════════════════════════════════════════════════════════

\begin{minipage}[t]{0.48\textwidth}
\begin{tcolorbox}[
    colback=gruppeA!10,
    colframe=gruppeA,
    title={\footnotesize\bfseries Kriterien Gruppe A (Partnerprüfung)},
    fonttitle=\bfseries,
    boxrule=0.5pt,
    top=2pt, bottom=2pt, left=4pt, right=4pt
]
\footnotesize
\begin{tabular}{@{}lr@{}}
1. ser/estar/hay korrekt & 30\% \\
2. Wortschatz Stadtviertel & 25\% \\
3. Grammatische Korrektheit & 20\% \\
4. Flüssigkeit + Interaktion & 25\% \\
\end{tabular}

\vspace{0.1cm}
\textit{Format: 3--4 Min/Paar, gegenseitig Fragen stellen}
\end{tcolorbox}
\end{minipage}
\hfill
\begin{minipage}[t]{0.48\textwidth}
\begin{tcolorbox}[
    colback=gruppeB!10,
    colframe=gruppeB,
    title={\footnotesize\bfseries Kriterien Gruppe B (Kurzvortrag)},
    fonttitle=\bfseries,
    boxrule=0.5pt,
    top=2pt, bottom=2pt, left=4pt, right=4pt
]
\footnotesize
\begin{tabular}{@{}lr@{}}
1. Inhalt + Begründung & 30\% \\
2. Fachvokabular Berufe & 25\% \\
3. Struktur + Kohärenz & 20\% \\
4. Flüssigkeit + Präsentation & 25\% \\
\end{tabular}

\vspace{0.1cm}
\textit{Format: 3--4 Min/Person, danach Fragen}
\end{tcolorbox}
\end{minipage}

\vspace{0.2cm}

% ═══════════════════════════════════════════════════════════════════════════
% SCHWIERIGKEITEN
% ═══════════════════════════════════════════════════════════════════════════

\begin{tcolorbox}[
    colback=white,
    colframe=spanisch,
    title={\footnotesize\bfseries Erwartete Schwierigkeiten},
    fonttitle=\bfseries,
    boxrule=0.5pt,
    top=2pt, bottom=2pt, left=4pt, right=4pt
]
\footnotesize
\begin{tabularx}{\textwidth}{@{}p{4cm}X@{}}
\textbf{ser/estar/hay Verwechslung} & SCAFFOLD-04-satzanfaenge nutzen, TAFELBILD-01 aufhängen \\
\textbf{E-Mail-Format unbekannt} & SCAFFOLD-04-beispiel-email zeigen (Muster) \\
\textbf{Unsicherheit bei Kriterien} & Konkrete Beispiele: ``Was bedeutet Flüssigkeit?'' \\
\textbf{A.A. schaut zu E.H.} & Für E-Mail trennen, eigenes Stadtviertel beschreiben lassen \\
\end{tabularx}
\end{tcolorbox}

\end{document}
